\documentclass[11pt,letterpaper]{article}

\usepackage[utf8]{inputenc}
\usepackage[T1]{fontenc}
\usepackage[spanish,es-noshorthands]{babel}
\usepackage[margin=2.4cm]{geometry}
\usepackage{amsmath,amssymb,bm}
\usepackage{booktabs}
\usepackage{longtable}
\usepackage{array}
\usepackage{listings}
\usepackage{xcolor}
\usepackage{enumitem}
\usepackage{hyperref}
\usepackage{siunitx}

\hypersetup{colorlinks=true, linkcolor=black, urlcolor=blue}
\sisetup{group-separator={,}, group-minimum-digits=4}

\lstset{
  basicstyle=\ttfamily\footnotesize,
  keywordstyle=\color{blue!60!black},
  commentstyle=\color{gray},
  stringstyle=\color{green!40!black},
  numbers=left, numberstyle=\tiny\color{gray}, numbersep=6pt,
  frame=single, rulecolor=\color{gray!40},
  breaklines=true, showstringspaces=false,
  language=Python, extendedchars=true, literate={í}{{\'i}}1 {á}{{\'a}}1 {é}{{\'e}}1 {ó}{{\'o}}1 {ú}{{\'u}}1 {ñ}{{\~n}}1
}

\title{\textbf{CC3104 --- Aprendizaje por Refuerzo}\\[0.3em]
       Examen Práctico: Auditoría del agente de gestión de inventario\\[0.5em]
       \large Grupo 4 --- La aproximación de función}
\author{%
\begin{tabular}{lcc}
\toprule
\textbf{Nombre} & \textbf{Carnet} & \textbf{Usuario Git} \\
\midrule
Edwin Jose Gabriel De Leon Garcia & 22809 & \texttt{EJGDLG} \\
Gustavo Adolfo Cruz Bardales      & 22779 & \texttt{G2309} \\
Josué Emanuel Say Garcia          & 22801 & \texttt{JosueSay} \\
Mathew Alexander Cordero Aquino   & 22982 & \texttt{donmatthiuz} \\
Carlos Alberto Valladares Guerra  & 221164 & \texttt{vgcarlol} \\
\bottomrule
\end{tabular}}
\date{\today}

\begin{document}
\maketitle

\tableofcontents
\bigskip

%==========================================================================
\section{Análisis del componente asignado}
%==========================================================================

\subsection*{Metodología de medición}
\addcontentsline{toc}{subsection}{Metodología de medición}

Para poder medir el error del aproximador se necesita un objetivo de referencia.
Como la función de transición del sistema es determinista, el MDP se resolvió de
forma exacta por iteración de valor con $\gamma = 0{,}99$, obteniendo $Q^{*}(s,a)$
para todos los estados alcanzables. Se usó la función de recompensa \emph{original}
del sistema entregado, no una corregida.

\begin{quote}
La especificación declara 200 estados (10 niveles de inventario $\times$ 5 de
vencimiento $\times$ 4 de demanda), pero \texttt{transition} devuelve
$\texttt{inventory} + \texttt{action} - \texttt{daily\_demand}$, que no es múltiplo
de 10 salvo cuando la demanda lo es. La cerradura del espacio de estados bajo la
función de transición contiene \textbf{4\,367 estados}, no 200. El agente pasa la
mayor parte del tiempo en estados que la especificación no declara.
\end{quote}

Todos los ajustes se hacen sobre los 4\,367 estados alcanzables, reservando 24
estados declarados como conjunto de prueba (Sección~\ref{sec:44}).

%--------------------------------------------------------------------------
\subsection{Entregable 4.1: el problema de escala}
%--------------------------------------------------------------------------

\subsubsection*{Qué componente domina el gradiente}

El aproximador es lineal, $Q(s,a;w) = w_a^{\top}\phi(s)$, y su actualización es la
regla LMS / semi-gradiente:
\[
  w_a \leftarrow w_a + \alpha\,\bigl[y - w_a^{\top}\phi(s)\bigr]\,
  \nabla_{w_a} Q(s,a;w),
  \qquad \nabla_{w_a} Q(s,a;w) = \phi(s).
\]
El gradiente \emph{es} el vector de características. Por lo tanto la magnitud del
paso que recibe la coordenada $i$ es proporcional a $\phi_i(s)$, y el aporte de esa
coordenada a la norma del gradiente es $\phi_i(s)^2 / \lVert\phi(s)\rVert^2$. Con
\begin{gather*}
  \phi(s) = \bigl[\,\texttt{inventory},\;\texttt{days\_to\_expiry},\;
  \texttt{demand\_idx}\,\bigr], \\[0.3em]
  \texttt{inventory} \in [0,100],\qquad
  \texttt{days} \in \{1,7,14,30,60\},\qquad
  \texttt{demand\_idx} \in \{0,1,2,3\},
\end{gather*}
el promedio sobre los 200 estados declarados da:

\begin{table}[h]
\centering
\begin{tabular}{lS[table-format=2.2]}
\toprule
\textbf{Componente} & {\textbf{Aporte a} $\lVert\phi\rVert^2$ \textbf{(\%)}} \\
\midrule
\texttt{inventory}       & 74.95 \\
\texttt{days\_to\_expiry} & 24.96 \\
\texttt{demand\_idx}      & 0.09 \\
\bottomrule
\end{tabular}
\caption{Participación de cada característica en la norma del gradiente.}
\end{table}

\noindent
\textbf{El nivel de inventario domina el gradiente con el 75\,\% y la demanda aporta
el 0{,}09\,\%}, unas 815 veces menos. La demanda determina si un producto se agota,
y es la componente que menos peso tiene en el gradiente. El peso asociado a la
demanda se mueve tres órdenes de magnitud más lento que el del inventario: para
cuando el aproximador ha aprendido algo sobre la demanda, el resto de los pesos ya
convergió a una solución que la ignora. El sesgo no es de capacidad sino de ritmo
de aprendizaje efectivo.

\subsubsection*{Divergencia del entrenamiento}

Para LMS con paso fijo $\alpha$, la condición de convergencia en media es
\[
  0 < \alpha < \frac{2}{\lVert\phi(s)\rVert^{2}}.
\]
El máximo observado es $\lVert\phi\rVert^{2}_{\max} = 11\,709$ (estado
$(90, 60, \texttt{crítico})$), de donde
\[
  \alpha_{\max} = \frac{2}{11\,709} = 1{,}708\times 10^{-4}.
\]
El código entregado usa $\alpha = 0{,}01$ por defecto en \texttt{update}, es decir
\textbf{58{,}5 veces por encima de la cota de estabilidad}. Se verificó
empíricamente ajustando $w$ por SGD sobre los estados de entrenamiento:

\begin{table}[h]
\centering
\begin{tabular}{lll}
\toprule
$\alpha$ & \textbf{Características originales} & \textbf{Características propuestas} \\
\midrule
$10^{-2}$   & diverge (\texttt{overflow}) & converge ($\lVert w\rVert_\infty \approx 3626$) \\
$10^{-3}$   & diverge (\texttt{overflow}) & --- \\
$1{,}4\times10^{-4}$ & converge & --- \\
$10^{-4}$   & converge & --- \\
$5\times10^{-2}$ & --- & converge ($\approx 3710$) \\
$10^{-1}$   & --- & converge ($\approx 3720$) \\
\bottomrule
\end{tabular}
\caption{Estabilidad empírica del ajuste por SGD. Con las características
originales, el $\alpha$ del código diverge; con las propuestas, $\alpha$ puede ser
hasta 100 veces mayor sin problemas.}
\end{table}

Adicionalmente, el número de condición de la matriz de Gram
$\frac{1}{N}\Phi^{\top}\Phi$ pasa de $\kappa = 1\,888$ con las características crudas
a $\kappa = 26{,}3$ con las mismas tres características normalizadas. Como la
velocidad de convergencia de un método de primer orden escala con $\kappa$, la
normalización sola acelera la convergencia en aproximadamente dos órdenes de
magnitud.

\subsubsection*{Defecto adicional: ausencia de intercepto}

$\phi$ no incluye un término constante y \texttt{weights} se inicializa en cero. Por
lo tanto, para el estado $(0, d, \texttt{bajo})$ se tiene $\phi = (0, d, 0)$, y para
$(0, \cdot, \texttt{bajo})$ con $d$ pequeño el vector es casi nulo: el modelo está
forzado a pasar por el origen. Como se muestra en la Sección~\ref{sec:44}, esto por
sí solo hace que el $R^2$ del aproximador original sea $-132{,}7$, o sea peor que
predecir constantemente el promedio de $Q^{*}$.

\subsubsection*{Normalización propuesta por componente}

Cada componente se lleva a $[0,1]$ usando una constante que proviene del propio
dominio, no de un escalado genérico:

\begin{table}[!ht]
\centering
\small
\begin{tabular}{p{3.1cm}p{4.6cm}p{7.2cm}}
\toprule
\textbf{Componente} & \textbf{Transformación} & \textbf{Justificación} \\
\midrule
\texttt{inventory} &
$\texttt{inv}/100 \in [0,1]$ &
100 es la capacidad real impuesta por el $\texttt{min}(100,\cdot)$ de
\texttt{transition}; no es una constante arbitraria. \\
\addlinespace
\texttt{days\_to\_expiry} &
$(d-1)/59 \in [0,1]$ más el indicador $\mathbf{1}[d \le 7]$ &
El rango declarado es $[1,60]$. El indicador replica la forma exacta de la
penalización $-10\cdot\mathbf{1}[d\le 7]$: sin él, ninguna combinación lineal de $d$
puede representar un escalón. \\
\addlinespace
\texttt{demand\_level} &
\emph{One-hot} con \texttt{bajo} como base &
El código ordinal $0,1,2,3$ impone dos supuestos falsos: que la demanda tiene orden
lineal en el valor y que el espaciado es uniforme. La demanda real es
$5,15,25,40$, cuyos incrementos son $10,10,15$. Además \texttt{bajo}$=0$ anula la
característica y hace indistinguibles dos estados con demanda baja. \\
\bottomrule
\end{tabular}
\caption{Normalización propuesta. Con ella $\lVert\phi\rVert^2 \le 5{,}93$ y
$\alpha_{\max} = 0{,}337$, frente a $1{,}708\times10^{-4}$ del original.}
\end{table}

%--------------------------------------------------------------------------
\subsection{Entregable 4.2: interacciones que el modelo lineal no captura}
%--------------------------------------------------------------------------

Un aproximador lineal sobre $\phi = [\texttt{inv}, d, \texttt{dem}]$ es aditivo: el
efecto marginal de cada variable es constante y no depende de las otras,
$\partial Q / \partial \texttt{inv} = w_{a,1}$ sin importar el valor de la demanda.
En gestión de inventario farmacéutico eso es falso. Las dos interacciones más
importantes son:

\paragraph{Interacción 1: inventario $\times$ demanda (riesgo de desabasto).}
El riesgo de \emph{stockout} no depende del inventario ni de la demanda por separado,
sino de su \textbf{cociente}: los días de cobertura
\[
  \texttt{cobertura}(s) = \frac{\texttt{inventory}}{\texttt{daily\_demand}}.
\]
Un inventario de 30 unidades es abundante con demanda baja (6 días de cobertura) y
crítico con demanda crítica (0{,}75 días). Un modelo aditivo asigna a
$\texttt{inv}=30$ el mismo aporte al valor en ambos casos. Un cociente no es
expresable como $w_1 x_1 + w_2 x_2$ para ningún $(w_1, w_2)$: es una imposibilidad
representacional, no un ajuste imperfecto. Se agrega
explícitamente como característica, saturada a 20 días y normalizada, junto con la
brecha inmediata $\max(0, \texttt{dem} - \texttt{inv})/40$, que marca los estados en
los que el desabasto ocurre en el paso siguiente.

\paragraph{Interacción 2: inventario $\times$ vencimiento $\times$ demanda (merma).}
Las unidades que van a vencerse son las que no alcanzarán a venderse antes del
vencimiento:
\[
  \texttt{merma}(s) = \max\bigl(0,\; \texttt{inventory} - \texttt{daily\_demand}
  \cdot \texttt{days\_to\_expiry}\bigr).
\]
Es un producto seguido de una rectificación. Los 41 productos vencidos por semana
en producción son exactamente esta cantidad, y el aproximador original no tiene
ninguna característica que se le parezca. Con
$d = 60$ y demanda crítica no hay merma con ningún inventario; con $d = 1$ y demanda
baja casi todo el inventario se pierde. Un término aditivo en $\texttt{inv}$ y otro
en $d$ no pueden distinguir esos dos casos. Se agrega también
$\texttt{inv}_n \cdot \mathbf{1}[d\le 7]$, que captura que el costo del inventario
alto se invierte de signo cuando el producto está por vencer.

%--------------------------------------------------------------------------
\subsection{Entregable 4.3: versión mejorada de \texttt{preprocess\_state}}
%--------------------------------------------------------------------------

El vector de características propuesto tiene 12 componentes:

\begin{table}[h]
\centering
\small
\begin{tabular}{clp{8.3cm}}
\toprule
\textbf{\#} & \textbf{Nombre} & \textbf{Definición y rol} \\
\midrule
0  & \texttt{bias}              & $1$ --- intercepto; permite representar estados con $\phi$ casi nulo. \\
1  & \texttt{inv\_n}            & $\texttt{inv}/100$ \\
2  & \texttt{inv\_n2}           & $(\texttt{inv}/100)^2$ --- no linealidad del costo de almacenamiento y saturación en la capacidad. \\
3  & \texttt{exp\_n}            & $(d-1)/59$ \\
4  & \texttt{near\_expiry}      & $\mathbf{1}[d \le 7]$ --- discontinuidad exacta de \texttt{reward}. \\
5--7 & \texttt{dem\_*}          & \emph{one-hot} de demanda (\texttt{medio}, \texttt{alto}, \texttt{crítico}; base \texttt{bajo}). \\
8  & \texttt{dias\_cobertura}   & $\min(\texttt{inv}/\texttt{dem},\,20)/20$ --- interacción 1. \\
9  & \texttt{riesgo\_merma}     & $\max(0, \texttt{inv} - \texttt{dem}\cdot d)/100$ --- interacción 2. \\
10 & \texttt{inv\_x\_near\_expiry} & $\texttt{inv}_n \cdot \mathbf{1}[d\le 7]$ \\
11 & \texttt{brecha\_stockout}  & $\max(0, \texttt{dem} - \texttt{inv})/40$ \\
\bottomrule
\end{tabular}
\caption{Características propuestas. Todas están en $[0,1]$, de modo que
$\lVert\phi\rVert^2 \le 12$ por construcción (máximo observado: $5{,}93$).}
\end{table}

\lstinputlisting[caption={\texttt{codigo/preprocesamiento\_mejorado.py}},
                 label={lst:pre}]{codigo/preprocesamiento_mejorado.py}

\subsubsection*{Cambios al aproximador}

\begin{enumerate}[leftmargin=*]
  \item \textbf{Paso normalizado.} \texttt{update} usa
  $\alpha_{\text{eff}} = \alpha / \lVert\phi(s)\rVert^{2}$ (regla LMS normalizada),
  que garantiza estabilidad para cualquier escala de entrada.
  \item \textbf{Inicialización optimista opcional.} Como ahora existe el intercepto,
  fijar $w_{a,0} = q_{\text{init}}$ produce $Q(s,a) \approx q_{\text{init}}$ para
  todo $s$, lo que da exploración dirigida sin tocar $\varepsilon$. Con el $\phi$
  original esto era imposible.
  \item \textbf{Compatibilidad.} La firma de \texttt{preprocess\_state} y de
  \texttt{predict}/\texttt{update} no cambia; solo cambia \texttt{n\_features} de 3
  a 12. El resto del sistema no requiere modificaciones.
\end{enumerate}

\subsubsection*{¿Es suficiente el aproximador lineal?}

Sí, para este dominio: $R^2 = 0{,}967$ sobre estados no vistos (entregable 4.4).
Las no linealidades de este MDP son conocidas y de forma cerrada --- un escalón en
$d = 7$, un producto rectificado para la merma, un cociente para la cobertura y una
saturación en la capacidad. Cuando la no linealidad se conoce, incluirla como
característica y ajustar un modelo lineal encima conserva convergencia garantizada
bajo aproximación lineal, pesos interpretables y costo de cómputo trivial.

Se requeriría aproximación no lineal si el estado se extendiera con variables cuya
interacción no se conoce de antemano: estacionalidad, sustitución entre productos o
demanda estocástica con correlación temporal.

El número de condición del bloque completo de 12 características es
$\kappa \approx 2{,}1\times10^{3}$, similar al del original, porque
\texttt{riesgo\_merma} es cero en la mayoría de los estados y queda casi colineal con
el resto. La normalización arregla la \emph{escala} (que es lo que causa la
divergencia), no elimina la colinealidad. Se mitiga con regularización ridge
($\lambda = 10^{-8}$ en los ajustes reportados) o eliminando la característica 10,
que es la de menor aporte marginal.

%--------------------------------------------------------------------------
\subsection{Entregable 4.4: comparación de error cuadrático medio}
\label{sec:44}
%--------------------------------------------------------------------------

\subsubsection*{Protocolo}

\begin{itemize}[leftmargin=*]
  \item Objetivo: $Q^{*}(s,a)$ exacto por iteración de valor ($\gamma = 0{,}99$,
        tolerancia $10^{-10}$) sobre los 4\,367 estados alcanzables.
  \item Conjunto de prueba: \textbf{24 estados representativos}, tomados de forma
        estratificada sobre los 200 estados declarados para cubrir todo el rango de
        inventario, los cinco niveles de vencimiento y los cuatro de demanda.
  \item Conjunto de entrenamiento: los 4\,343 estados restantes. Los estados de
        prueba \emph{no} se usan para ajustar.
  \item Ajuste: mínimos cuadrados exacto por acción (ecuaciones normales con ridge
        $\lambda = 10^{-8}$). Se usa la solución óptima en lugar de SGD para que la
        comparación mida \textbf{capacidad de representación} y no calidad del
        optimizador; el efecto sobre el optimizador se reporta por separado en 4.1.
  \item Métrica: $\mathrm{MSE} = \frac{1}{|S|\cdot|A|}\sum_{s}\sum_{a}
        \bigl(Q^{*}(s,a) - w_a^{\top}\phi(s)\bigr)^{2}$.
\end{itemize}

\subsubsection*{Resultado global}

\begin{table}[h]
\centering
\begin{tabular}{lS[table-format=7.1]S[table-format=7.1]S[table-format=-3.3]}
\toprule
\textbf{Aproximador} & {\textbf{MSE entren.}} & {\textbf{MSE prueba}} & {$\bm{R^2}$ \textbf{prueba}} \\
\midrule
Original (3 caract., sin intercepto)      & 1635206.0 & 2759751.7 & -132.716 \\
Original + intercepto (4 caract.)         & 1041.6    & 990.1     & 0.952 \\
\textbf{Propuesto (12 caract.)}           & 318.2     & 678.4     & 0.967 \\
\bottomrule
\end{tabular}
\caption{Comparación de MSE. La varianza de $Q^{*}$ es
$\sigma^2 = 20\,638{,}8$; $R^2 = 1 - \mathrm{MSE}/\sigma^2$.
Reducción del MSE de prueba respecto al original: \textbf{99{,}98\,\%}.}
\label{tab:mse}
\end{table}

La fila intermedia (características originales más un intercepto) es una ablación.
El salto de $2{,}76\times10^{6}$ a $990$ se debe casi por completo al intercepto
faltante, no a las características nuevas. Aislado ese efecto, las interacciones del
entregable 4.2 aportan una \textbf{reducción adicional del 31{,}5\,\%} del MSE de
prueba ($990 \to 678$) y suben el $R^2$ de $0{,}952$ a $0{,}967$.

El ajuste original es \emph{peor que la predicción constante} ($R^2 < 0$). Un
aproximador con $R^2$ negativo alimentando una política greedy
$\arg\max_a Q(s,a)$ no ordena las acciones mejor que el azar en los estados donde
$\phi$ es pequeño, que son los de inventario bajo, donde ocurren los desabastos.

\subsubsection*{Detalle por estado (24 estados de prueba)}

\begin{longtable}{rrlS[table-format=4.1]cS[table-format=8.1]S[table-format=4.1]}
\toprule
\textbf{Inv.} & $\bm{d}$ & \textbf{Demanda} & {$\bm{\max_a Q^{*}}$} & $\bm{a^{*}}$ &
{\textbf{MSE original}} & {\textbf{MSE propuesto}} \\
\midrule
\endfirsthead
\toprule
\textbf{Inv.} & $\bm{d}$ & \textbf{Demanda} & {$\bm{\max_a Q^{*}}$} & $\bm{a^{*}}$ &
{\textbf{MSE original}} & {\textbf{MSE propuesto}} \\
\midrule
\endhead
\bottomrule
\endfoot
  0 &  1 & bajo    & 3767.5 & 50 & 13610118.3 &  233.3 \\
  0 & 14 & bajo    & 3826.1 & 50 &  9319062.9 &  216.8 \\
  0 & 60 & bajo    & 4174.6 & 50 &   887390.8 &  653.9 \\
 10 &  7 & medio   & 3762.6 & 50 &  6820040.4 &  327.4 \\
 10 & 30 & medio   & 3961.0 & 50 &  2497293.7 &  211.5 \\
 20 &  1 & medio   & 3772.5 & 50 &  7258968.9 &  506.5 \\
 20 & 14 & alto    & 3813.7 & 50 &  2344597.2 &  105.3 \\
 20 & 60 & alto    & 4162.2 & 50 &   335822.6 &  457.4 \\
 30 &  7 & alto    & 3767.6 & 50 &  2600258.0 &  982.1 \\
 30 & 30 & crítico & 3895.1 & 50 &     2425.2 & 1067.7 \\
 40 &  1 & crítico & 3726.0 & 50 &  1196949.6 &  264.0 \\
 40 & 14 & crítico & 3784.5 & 50 &   210012.0 & 1988.1 \\
 50 &  1 & bajo    & 3797.5 & 50 &  6165322.4 &  711.7 \\
 50 & 14 & bajo    & 3856.0 & 50 &  3409680.5 &    2.9 \\
 50 & 60 & bajo    & 4204.5 & 50 &    69797.3 &   94.4 \\
 60 &  7 & medio   & 3797.5 & 50 &  1998061.6 &  293.5 \\
 60 & 30 & medio   & 3995.9 & 50 &   146148.3 &  140.5 \\
 70 &  1 & medio   & 3800.0 & 50 &  2222812.5 &  338.2 \\
 70 & 14 & alto    & 3856.0 & 50 &   121051.8 &   31.5 \\
 70 & 60 & alto    & 4204.5 & 50 &  3107585.7 &  281.0 \\
 80 &  7 & alto    & 3800.0 & 50 &   177655.3 &  503.7 \\
 80 & 30 & crítico & 3993.4 & 50 &  1300516.0 & 1923.6 \\
 90 &  1 & crítico & 3800.0 & 50 &      444.9 & 3752.1 \\
 90 & 14 & crítico & 3858.5 & 50 &   432025.0 & 1194.6 \\
\caption{MSE por estado, promediado sobre las 6 acciones. El aproximador propuesto
gana en 23 de los 24 estados.}
\label{tab:detalle}
\end{longtable}

\subsubsection*{Lectura de la tabla}

\begin{itemize}[leftmargin=*]
  \item El aproximador propuesto gana en \textbf{23 de 24} estados. El único
        estado donde pierde es $(90,1,\text{crítico})$, con inventario máximo y
        vencimiento inmediato: ahí $\lVert\phi_{\text{orig}}\rVert$ es grande y el
        modelo sin intercepto acierta la escala de $Q^{*}$ por coincidencia de
        magnitud.
  \item El error del original crece a medida que el inventario baja
        ($1{,}4\times10^{7}$ en $\texttt{inv}=0$), exactamente el efecto del
        intercepto ausente descrito en 4.1.
  \item El error residual del modelo propuesto se concentra en los estados de
        demanda crítica con vencimiento cercano
        ($(90,1,\text{crítico})$: 3752). Son los estados de frontera del dominio y
        los peor cubiertos por el conjunto de entrenamiento.
  \item La acción óptima es \texttt{pedir 50} en los 24 estados de prueba. Con la
        función de recompensa original, pedir el máximo \emph{sí} es la política
        óptima del MDP: el ``94\,\% de estados con \texttt{pedir\_50}'' de las
        métricas de entrenamiento no es un fallo de aprendizaje, sino el resultado
        de una recompensa mal especificada. Ninguna corrección del preprocesamiento
        lo modifica.
\end{itemize}

\subsubsection*{Reproducibilidad}

\begin{lstlisting}[language=bash, numbers=none, caption={Reproducción de los resultados}]
python3 codigo/experimento_mse.py
\end{lstlisting}

\noindent
No requiere dependencias externas (Python puro). Los valores crudos, incluidos los
pesos ajustados, quedan en \texttt{resultados/resultados\_44.json}.

\clearpage
%==========================================================================
\section{Respuestas a las preguntas de integración}
%==========================================================================

\subsection{Pregunta 1: Diagnóstico sistémico}

Con base en el MDP estocástico formulado por el Grupo~1, la función de recompensa multicomponente diseñada por el Grupo~2, la ablación empírica y curvas proyectadas por el Grupo~6, y el análisis de aproximación y estabilidad de nuestro Grupo~4, \textbf{predecimos categóricamente que el agente NO aprenderá una política mejor si únicamente se aplican las correcciones del MDP y de la función de recompensa sin modificar el algoritmo de aprendizaje ni la estrategia de exploración}.

Aunque la corrección conjunta del MDP y de la señal de recompensa repara la especificación matemática del entorno y reorienta teóricamente los incentivos hacia políticas de inventario sostenibles, el proceso de optimización del agente colapsa inevitablemente debido a tres fallas estructurales y cuantitativas:

\begin{enumerate}
    \item \textbf{Eficacia teórica aislada de la recompensa de G2:}
    Bajo la función corregida del Grupo~2 ($R_{\text{corregida}}$), en una simulación de $1\,000$ episodios y $30$ pasos por episodio (semilla $3104$), la política viciada de ``pedir siempre $50$'' es severamente penalizada, acumulando una recompensa media de $-842.29$ ($\sigma = 476.22$), con un inventario promedio saturado de $97.98$ unidades y $18.11$ pasos por episodio en ventana de vencimiento. Por el contrario, una política balanceada alcanza una recompensa acumulada media positiva de $+144.40$ ($\sigma = 67.81$), reduciendo el inventario promedio a $5.76$ unidades, los pasos en riesgo de vencimiento a $7.48$ y manteniendo $0.00$ pasos con desabasto (\textit{stockouts}). Esto evidencia que la señal de recompensa corregida contiene el gradiente de incentivos adecuado para guiar al agente.

    \item \textbf{Explosión del espacio de búsqueda de G1 frente a la inestabilidad de G3 y G4:}
    Para restaurar la propiedad de Markov y dotar de dinamismo temporal a la demanda, el Grupo~1 amplió la tupla de estado a 5 variables: $s_t = (I_t, E_t, L_t, T_t, D_t)$, incorporando las unidades en tránsito ($T_t \in \{0, 10, 20, 30, 40, 50\}$) y la tendencia de demanda ($D_t \in \{\text{bajando, estable, subiendo}\}$). Dicha corrección expande el espacio de estados de $200$ nominales a $3\,600$ estados discretos, y el espacio de pares estado-acción de $1\,200$ a $21\,600$ (un factor de multiplicación exacto de $18\times$).

    Si se mantiene el algoritmo de aprendizaje original:
    \begin{itemize}
        \item El valor fijo $\alpha = 0.9$ es catastrófico para la convergencia. Como demostró analíticamente nuestro Grupo~4 mediante la condición LMS ($0 < \alpha < \frac{2}{\|\phi(s)\|^2}$), el máximo paso admisible antes de divergir numéricamente es $\alpha_{\max} \approx 1.708 \times 10^{-4}$. Un valor $\alpha = 0.9$ se encuentra tres órdenes de magnitud por encima de la cota de estabilidad.
        \item Al contar con un espacio de búsqueda $18$ veces mayor, cada par $(s, a)$ se visita con una frecuencia mucho menor durante $1\,000$ episodios, haciendo que las actualizaciones con un $\alpha$ constante y excesivo generen oscilaciones caóticas y divergencia en las aproximaciones de valor en lugar de un aprendizaje asintótico.
    \end{itemize}

    \item \textbf{Evidencia empírica de ablación (G6) y deficiencia de cobertura ($\epsilon = 0.05$):}
    El experimento de ablación controlado del Grupo~6 valida de forma empírica la incapacidad del optimizador original para alcanzar la política esperada:
    \begin{itemize}
        \item Al ejecutar el entrenamiento aplicando \textbf{únicamente la recompensa corregida de G2} pero manteniendo el algoritmo original ($\alpha = 0.9$, $\epsilon = 0.05$), el agente no logró acercarse a la cota óptima ($+144.40$): la recompensa acumulada colapsó a un promedio en cola de $-627.73$, con un error TD absoluto terminal inestable de $35.80$ y una varianza masiva de $39\,849.11$.
        \item La política colapsó de manera prematura a una regla subóptima y pasiva: la acción dominante pasó a ser pedir $10$ unidades en apenas el $25.5\%$ de los estados visitados, dejando al resto del espacio sin converger adecuadamente.
        \item Bajo $\epsilon = 0.05$, el agente actúa de manera puramente \textit{greedy} el $95\%$ del tiempo. Con un horizonte de $30$ pasos por episodio, el Grupo~1 demostró que se requerirían aproximadamente $18\,000$ episodios (en lugar de los $1\,000$ asignados) para conservar una densidad de visitas equiparable por estado. Con una tasa estática tan baja, los estados críticos y de demanda alta rara vez se visitan, perpetuando desabastos en la inferencia en producción.
    \end{itemize}
\end{enumerate}

Corregir el MDP y la recompensa repara exclusivamente el modelado del entorno y la métrica de desempeño (el ``qué'' optimizar), pero deja inalterada la dinámica del algoritmo y la exploración (el ``cómo'' aprender). Sin un algoritmo numéricamente estable (mediante el decaimiento de $\alpha$ propuesto por el Grupo~3 o el paso normalizado LMS de nuestro Grupo~4) y sin un mecanismo de exploración eficiente en espacios de alta dimensionalidad (Grupo~5), el agente es incapaz de converger y termina atrapado en políticas degeneradas.

\vspace{4cm}

\subsection{Pregunta 2: Causa raíz}

\vspace{4cm}

\subsection{Pregunta 3: Plan de corrección mínimo viable}

\vspace{4cm}

\subsection{Pregunta 4: Compatibilidad de correcciones}

\vspace{4cm}

\clearpage
%==========================================================================
\section{Reflexión grupal}
%==========================================================================

Durante la fase previa de análisis individual, nuestro equipo llegó a la sesión con una hipótesis clara: el mal desempeño del agente se debía principalmente a la deficiente ingeniería de características y a la escala sin normalizar en el preprocesamiento. Al constatar que el inventario concentraba el 75\,\% de la norma del gradiente frente a apenas el 0.09\,\% de la demanda, y al comprobar analíticamente que la tasa $\alpha = 0.9$ violaba la condición de estabilidad LMS ($\alpha_{\max} \approx 1.708 \times 10^{-4}$), asumíamos que corrigiendo la representación lineal e incorporando las interacciones no lineales resolveríamos el problema. Tras verificar que nuestro modelo propuesto reducía el MSE de prueba en un 99.98\,\%, considerábamos que el cuello de botella estaba resuelto.

El contraste de resultados en la sesión presencial cambió por completo este diagnóstico. Al revisar la función de recompensa con el Grupo~2, descubrimos que el entorno sufría de un severo \textit{reward hacking} que premiaba pedir siempre 50 unidades; por tanto, aunque nuestro aproximador fuera perfecto con un MSE nulo sobre los 4\,367 estados alcanzables, el agente seguiría aprendiendo una política destructiva. Asimismo, el Grupo~1 nos demostró que el estado original carecía de la propiedad de Markov y que, al corregirlo con el inventario en tránsito y la tendencia de demanda, el espacio de estados se multiplicaba por 18 ($3\,600$ estados), obligando a que cualquier esquema de aproximación deba gestionar una dimensionalidad mucho mayor. Finalmente, los experimentos de ablación del Grupo~6 confirmaron que aun aplicando la recompensa corregida, el agente colapsaba a recompensas negativas ($-627.73$) si no se estabilizaban el algoritmo y la aproximación de forma conjunta.

La lección fundamental de esta auditoría es que en Aprendizaje por Refuerzo no es posible resolver un problema optimizando componentes de forma aislada. La aproximación de funciones no opera en el vacío, sino como el puente que traduce el estado del entorno (MDP) y transmite el gradiente de la señal de desempeño (recompensa). De nada sirve perfeccionar la representación de características si la recompensa premia una conducta errónea o si el entorno oculta variables críticas para la toma de decisiones.

\clearpage
\appendix
%==========================================================================
\section{Anexo: uso de IA generativa}
%==========================================================================

Se usó IA generativa como apoyo de herramienta, no para el análisis del componente.
A continuación los prompts empleados y la respuesta que se obtuvo de cada uno.

\subsection*{Prompt 1: compilación del documento}

\begin{quote}
\small
\ttfamily
No tengo \texttt{pdflatex} instalado en el sistema y estoy en NixOS. Necesito compilar
\texttt{informe\_grupo4.tex} a PDF sin instalar nada de forma permanente. El documento usa
\texttt{babel} en español, \texttt{booktabs}, \texttt{longtable}, \texttt{siunitx},
\texttt{enumitem} y \texttt{listings} con \texttt{lstinputlisting}. Dame el comando exacto
con \texttt{nix-shell -p}, indicando qué esquema de TeX Live necesito para que no
falte ningún paquete, y explícame por qué hay que correr \texttt{pdflatex} más de una vez.
\end{quote}

\noindent
\textbf{Respuesta obtenida:}

\begin{lstlisting}[language=bash, numbers=none]
nix-shell -p texlive.combined.scheme-full --run \
  'pdflatex -interaction=nonstopmode -halt-on-error informe_grupo4.tex && \
   pdflatex -interaction=nonstopmode -halt-on-error informe_grupo4.tex'
\end{lstlisting}

\subsection*{Prompt 2: escritura de las fórmulas}

\begin{quote}
\small
\ttfamily
Convierte estas expresiones a \LaTeX{} con \texttt{amsmath}, en modo display y
tipográficamente correctas:
(1) la regla de actualización semi-gradiente para un aproximador lineal
$Q(s,a;w)=w_a^\top \phi(s)$, dejando explícito que el gradiente respecto a $w_a$ es
$\phi(s)$;
(2) la condición de estabilidad de LMS $0 < \alpha < 2/\lVert\phi(s)\rVert^2$;
(3) el error cuadrático medio promediado sobre estados y acciones.
Requisitos: usa \texttt{\textbackslash lVert}/\texttt{\textbackslash rVert} y no
\texttt{||}; nombres de variables del código en \texttt{\textbackslash texttt}
dentro de modo matemático; y si una línea se pasa del margen, parte la ecuación con
\texttt{gather*} en vez de dejar overfull hbox.
\end{quote}

\noindent
\textbf{Respuesta obtenida:}

Regla de actualización semi-gradiente:
\[
  w_a \leftarrow w_a + \alpha\,\bigl[y - w_a^{\top}\phi(s)\bigr]\,
  \nabla_{w_a} Q(s,a;w),
  \qquad \nabla_{w_a} Q(s,a;w) = \phi(s).
\]

\noindent
Condición de estabilidad de LMS:
\[
  0 < \alpha < \frac{2}{\lVert\phi(s)\rVert^{2}}.
\]

\noindent
Error cuadrático medio sobre estados y acciones:
\[
  \mathrm{MSE} = \frac{1}{|S|\cdot|A|}\sum_{s}\sum_{a}
  \bigl(Q^{*}(s,a) - w_a^{\top}\phi(s)\bigr)^{2}.
\]

\end{document}
